% docs/common/latex/conventions.tex - front matter shared by the three books (PDF only): how the books are typeset.
% Hand-written LaTeX (no Markdown counterpart); \input by the main .tex of each book.
%   Compile: cmake --build <build> --target pdf, or python3 -I tools/build_docs_pdf.py --source-root . --build-dir DIR
%   (figures, then xelatex passes + makeindex -s docs/common/latex/simcpp.ist); by hand see docs/README.md.
%   XeLaTeX; class docs/common/latex/simcpp.cls (TeX Live packages listed there); fonts TeX Gyre Pagella,
%   TeX Gyre Heros, DejaVu Sans Mono. The figures are built first from docs/{manual,translation}/figures/*.tex
%   (preview.sty) into PDFs.
\chapter*{Conventions}\addcontentsline{toc}{chapter}{Conventions}\markboth{Conventions}{}\label{conventions}
\simcabstract{How to read the typography, the boxes and the paths of this book.}

\section*{Typefaces}
Body text is set in a serif face; headings, tables and figures in a sans-serif face. Everything that is typed or read
by a program is set in a monospaced face and coloured: file names (\code{data/matrices}), input parameters
(\code{doing\_pion}), command-line options (\code{-{}-threads}), classes and functions (\code{SimcRun::step}). In the
user manual an input parameter is a link to its entry in the parameter reference. Links inside the book are blue;
a reference to a document outside the book is given in a footnote with its path in the repository.

\section*{Boxes}
Command lines are shown in a shaded box with a blue bar; they are typed in a shell:
\begin{simcshell}[basicstyle=\ttfamily\footnotesize]
simc++ $EX/h_elastic_rad.inp -o out --threads 4
\end{simcshell}
Source code and configuration files (C++, JSON, CMake) are framed by thin rules and coloured:
\begin{simccode}[basicstyle=\ttfamily\footnotesize, style=simcjson]
{"pipeline": {"rng": {"type": "ranlux"}}}
\end{simccode}
Program output, file contents and text diagrams have a light grey frame. A long line continues on the next line
after the sign \textcolor{simcblue}{$\hookrightarrow$}; a block copied from a source file names that file above it.

\begin{simcnote}
A note adds information that helps to understand the text.
\end{simcnote}
\begin{simctip}
A tip shows a shorter or better way to do something.
\end{simctip}
\begin{simcwarning}
A warning points at a mistake that is easy to make and costly.
\end{simcwarning}
\begin{simclegacy}
A ``Legacy SIMC'' box explains a behaviour inherited from the Fortran SIMC.
\end{simclegacy}

\section*{Paths}
The commands use these locations; adapt them to your machine. Relative paths such as \code{src/driver} are relative to
the source tree.

\begin{simctable}{small}{2}{>{\raggedright\arraybackslash}p{0.25\linewidth}>{\raggedright\arraybackslash}p{0.71\linewidth}}
\toprule\rowcolor{simchead}\simcth{variable} & \simcth{meaning}\\\midrule\endhead
\bottomrule\endlastfoot
\code{\$SIMCPP\_SRC} & the git checkout of SIMC++ (the source tree)\\
\code{\$SIMCPP\_BUILD} & the CMake build directory\\
\code{\$SIMCPP\_PREFIX} & the installation prefix (\code{bin/}, \code{lib/}, \code{include/}, \code{share/simc/data})\\
\code{\$SIMC\_FORTRAN} & optional: the legacy Fortran SIMC, used only by validation tests\\
\code{\$SIMC\_REFERENCE} & optional: reference runs and tables of the legacy program, used only by validation tests\\
\code{\$EX} & the example inputs of the user manual (\code{\$SIMCPP\_PREFIX/share/simc/docs/manual/markdown/user/examples})\\
\end{simctable}

\section*{Units}
Unless stated otherwise: energies and momenta in MeV and MeV/c, lengths in cm, angles in rad (inputs in degrees
where the input says so), \code{delta} in percent, slopes as dx/dz and dy/dz.

\section*{Index}
The index at the end lists the input parameters, the command-line options and the classes and interfaces of
SIMC++ under their own headings, with the pages where they are described or used.
